% Passes 9725--9788: transporter sign obstruction, full unitary glue symmetry, coherent readout, local-field bridge.
\subsubsection*{Passes 9725--9788: the transporter obstruction is Witt sign, and the $F_9$ package has a local-field lift}
The natural point-scheme version of the Hall--Janko/$Q^-(5,3)$ bridge is now
excluded.  The $252$ nonsingular $Q^-$ points split into two norm fibers of size
$126$, and the basic $O^-(6,3)$ relation valencies are $45,80,36,90$.  Hence
all invariant undirected union degrees lie in
\[
\{0,36,45,80,81,90,116,125,126,135,161,170,171,206,215,251\},
\]
which excludes the $G_2(4)$ fixed-edge $D$-cell degree $66$.  Likewise a
single $126$-point norm fiber has only degrees $45,80,125$, excluding the
$G_2(4)$ $B\cup C$ degree $31$.  Any surviving bridge must therefore use a
derived scheme such as two-spaces or Lagrangians rather than point orthogonality.

The intrinsic symmetry of the transverse Golay/$E_6$ glue pair is much larger
than its signed-coordinate symmetry.  Inside
\[
C_{\operatorname{Sp}(12,3)}(R)\cong U(6,3),
\]
the complete ordered-pair stabilizer is
\[
\boxed{O^-(6,3)},\qquad |O^-(6,3)|=26,127,360,
\]
with projective order $13,063,680$.  Allowing the two glue halves to exchange
doubles the stabilizer to $52,254,720$.  Thus the earlier $\{\pm I\}$ result
is correctly interpreted as signed-monomial rigidity, not absolute unitary
rigidity.

The $A_2$ line-orientation quotient and the $G_2(4)$ edge-orientation quotient
are both $C_2$: the former is
$3^3{:}S_4\to C_2$ with kernel $3^3{:}A_4$, while the latter is endpoint
reversal, exchanging the unique nonsymmetric rank-$14$ pair
$AD\leftrightarrow BB/CC$.  Once an external bridge matches the underlying
unoriented objects, the orientation identification is forced because
$\operatorname{Aut}(C_2)$ is trivial.  However the parent stabilizers of orders
$648$ and $24,192$ cannot embed into one another, so the bit itself is not the
transporter.

More strongly, the presently distinguished transporter is impossible.  On the
glue side the exchanger $R$ satisfies $R^2=-I$ and the canonical Lagrangian
half carries $Q^-(5,3)$ with $112$ singular projective points.  On the exact
ATLAS Suzuki side the exchanger $x$ is symplectic with $x^2=-I$, but its
canonical M12 polarization half carries $Q^+(5,3)$ with $130$ singular points.
If a symplectic map $T$ both conjugated $R$ to $x$ and mapped one distinguished
half to the other, then the induced forms
\[
K(u,Rv)\quad\text{and}\quad J(Tu,xTv)
\]
would be congruent, preserving Witt sign, which would force $112=130$.
Therefore
\[
\boxed{\text{no canonical-half-preserving }R\text{ transporter exists}.}
\]
The obstruction is selective rather than fatal.  For fixed $R$, the unitary
centralizer has two orbits of transverse Lagrangians:
\[
\begin{array}{c|c|c}
\text{half-form}&\text{stabilizer}&\text{orbit size}\\\hline
Q^+(5,3)&O^+(6,3)&7,530,558,336\\
Q^-(5,3)&O^-(6,3)&6,992,661,312.
\end{array}
\]
Thus the same glue-derived $R$ admits billions of plus-type alternative
polarizations capable of clearing the Suzuki Witt-sign obstruction.  The next
problem is to select one canonically.

A coherent finite readout model was also tested on $40\times3$ complex
amplitudes: forty $W33$ ports and three internal qutrit/OAM--time-bin modes,
with unitary propagation $\exp(-i\theta A_{W33}/12)$, bounded phase disorder,
loss, random internal $3\times3$ unitary mixing, finite detector shots and a
noisy dark monitor.  A line-plus-dark classifier remained correct in all
$1,500$ trials under both the moderate and heavier seeded stress profiles.  A
conditional direct phase tag for the $A_2$ orientation bit survived the
moderate profile but failed under heavier phase/loss disorder before carrier
identity failed.  This is a robustness simulation, not a measured hardware
Hamiltonian.

Two cross-lane consequences are especially suggestive.  First, using the
corrected Leech result that the canonical $V_2$ generator contains $4,095$
intrinsic type-$8$ frames, together with the $416$ Hall--Janko vertices and
$20,800$ $G_2(4)$ edges,
\[
\gcd(4095,416,20800)=13.
\]
Hence any common selector with constant fiber size on all three canonical sets
can have only one nontrivial cardinality,
\[
\boxed{13=\Phi_3(3)}.
\]
Existence is not asserted; the parallel frame reconstruction additionally
shows that naive mod-$2$ frame coordinates cannot construct it.

Second, the $F_9$ glue phase space and the parallel Bruhat--Tits cyclotomic
filtration combine naturally in local-field arithmetic.  The relation
$R^2=-I$ supplies the unramified quadratic residue extension
$\mathbb F_9/\mathbb F_3$, lifted by $\mathbb Q_3(i)/\mathbb Q_3$, while the
order-$9$ cyclotomic tower supplies the totally ramified degree-$6$ extension
$\mathbb Q_3(\zeta_9)$ with uniformizer $\pi=1-\zeta_9$.  Their compositum
\[
L=\mathbb Q_3(i,\zeta_9)
\]
has
\[
(e,f)=(6,2),\qquad [L:\mathbb Q_3]=12,
\]
residue field $\mathbb F_9$, and $3=u\pi^6$.  Consequently
$\mathcal O_L/3\mathcal O_L$ has six successive $\pi$-adic graded pieces, each
$\mathbb F_9$, so its associated graded additive space is
\[
\boxed{\mathbb F_9^6\cong\mathbb F_3^{12}}.
\]
This gives an arithmetic lift candidate for the independently discovered glue
phase space: $R$ is the unramified direction and $I-M$ the ramified uniformizer
direction.  No integral Niemeier--$\mathcal O_L$ identification is claimed.
